% The order of a 5-chromatic unit-distance graph is at least 29.
%
% Compiles with pdflatex or tectonic, TeX Live 2023 or later.  No BibTeX: the
% bibliography is a manual thebibliography.  No figures, no packages outside the
% TeX Live base.  Single file.
\documentclass[11pt,reqno]{amsart}

\usepackage[T1]{fontenc}
\usepackage[utf8]{inputenc}
\usepackage{amsmath}
\usepackage{amssymb}
\usepackage{amsthm}
\usepackage[hidelinks]{hyperref}

\theoremstyle{plain}
\newtheorem{theorem}{Theorem}
\newtheorem{proposition}[theorem]{Proposition}

\theoremstyle{definition}
\newtheorem{remark}[theorem]{Remark}

\newcommand{\R}{\mathbb{R}}

\begin{document}

\title[A 5-chromatic unit-distance graph has at least 29 vertices]
      {The order of a 5-chromatic\\ unit-distance graph is at least $29$}

\author{Andrei-C\u{a}lin Nedelcu}
% TODO before submission: amsart expects an address and AMS journals require
% one.  Replace the line below with your institution, or with a plain postal
% town if you are submitting unaffiliated -- both are acceptable.
\address{\textbf{[AFFILIATION TO BE SUPPLIED]}}
\email{calin.nedelcu08@gmail.com}

\subjclass[2020]{05C15, 52C10, 05C62}
\keywords{Hadwiger--Nelson problem, unit-distance graph, chromatic number of the
plane, Erd\H{o}s unit distance problem}

\begin{abstract}
Let $v_5$ be the least number of vertices of a 5-chromatic unit-distance graph in
the Euclidean plane. The best published bound is $v_5 \ge 28$, due to de Grey and
Parts. We observe that their own edge bound, combined with the unit-distance
counts of Alexeev, Mixon and Parshall published two years later, gives
$v_5 \ge 29$ with no new computation. The argument is two lines; what it needs is
a table that did not exist when the vertex bound was derived. We also record
exactly what $30$ and $31$ would require, and note that one of the two obstacles
de Grey and Parts name for their own method is removable for polygonal tilings.
\end{abstract}

\maketitle

\section{The two ingredients}

Write $u(n)$ for the maximum number of unit distances determined by $n$ points in
the plane, and $e_5$ for the least number of edges of a 5-chromatic
unit-distance graph in $\R^2$.

\begin{theorem}[de Grey and Parts \cite{dGP}]\label{thm:A}
$e_5 \ge 99$ and $v_5 \ge 28$.
\end{theorem}

Their bound on $e_5$ comes from a refinement of Pritikin's uncovered-density
scheme together with the Gwyn--Stavrianos badness argument: if $p_4$ is the
probability that a random unit segment is monochromatic under the best
$4$-colouring of the plane available to the method, then $e_5 \ge \lceil 1/p_4
\rceil$. They state the two bounds together, in the form ``we get the bounds:
$e_5 \ge 99$, $v_5 \ge 28$, $e_6 \ge 182$, $v_6 \ge 42$''.

\begin{theorem}[Alexeev, Mixon and Parshall \cite{AMP}]\label{thm:B}
For $22 \le n \le 30$,
\[
\begin{array}{l|ccccccccc}
 n            & 22 & 23 & 24 & 25 & 26 & 27 & 28 & 29 & 30 \\ \hline
 u(n) \ge     & 60 & 64 & 68 & 72 & 76 & 81 & 85 & 89 & 93 \\
 u(n) \le     & 61 & 66 & 72 & 78 & 84 & 90 & 96 & 103 & 110
\end{array}
\]
\end{theorem}

The lower bounds are earlier constructions; the upper bounds are that paper's
contribution.

\section{The observation}

\begin{proposition}\label{prop:29}
$v_5 \ge 29$.
\end{proposition}

\begin{proof}
Suppose $G$ is a 5-chromatic unit-distance graph in the plane on $n \le 28$
vertices. Its edges are unit distances among $n$ points, so $\lvert E(G)\rvert
\le u(n)$. The function $u$ is non-decreasing, since adding a point removes no
unit distance, so $\lvert E(G)\rvert \le u(28) \le 96$ by Theorem~\ref{thm:B}.
But $\lvert E(G)\rvert \ge e_5 \ge 99$ by Theorem~\ref{thm:A}, and $96 < 99$.
\end{proof}

\begin{remark}
The bound is exactly $29$ by this route and no further: $u(29) \le 103$, which
does not contradict $e_5 \ge 99$.
\end{remark}

\section{Why it was available but not taken}

De Grey and Parts convert edges to vertices using the bounds of \'Agoston and
P\'alv\"olgyi, which in the relevant range allow at most $97$ edges on $27$
vertices and at most $102$ on $28$. So $27$ was excluded and $28$ was not: their
bound is exactly what their inputs allow. Their paper is from March 2023;
\cite{AMP} is from December 2024, revised February 2025, and is strictly better
at every $n$ in this range. At $n = 28$ it replaces $102$ by $96$, and $28$ no
longer survives.

Nothing else is needed, and nothing else is claimed. This is an observation about
two existing results, not a new method.

\section{What $30$ and $31$ require}

Two independent levers, both engineering rather than new theory.

\subsection*{Lever 1: raise $e_5$}
Against the table of Theorem~\ref{thm:B}, held fixed,
\[
  e_5 \ge 104 \implies u(29) \le 103 < 104 \implies v_5 \ge 30,
\]
\[
  e_5 \ge 111 \implies u(30) \le 110 < 111 \implies v_5 \ge 31.
\]
Reaching $e_5 \ge 104$ needs $p_4 < 1/104 = 0.0096154$, about $5\%$ below the
value behind $99$. De Grey and Parts name their own obstacle: they compute $p_4$
by numerical integration and write that they ``encountered some problems with the
numerical integration function \texttt{NIntegrate} in the general case, and were
forced to limit ourselves to relatively simple tilings'', nominating tilings with
wavy borders as the promising family.

For a \emph{polygonal} tiling that obstacle is removable. The integrand
$M_k(\sigma,\varphi)$ of
\[
  p_k = \frac{1}{2\pi S}\int_S \mathrm{d}\sigma \int_0^{2\pi} \mathrm{d}\varphi \;
        M_k(\sigma,\varphi)
\]
is piecewise constant: a unit segment's two endpoints each lie in some cell, and
$M_k$ depends only on which pair of cells. So $p_k$ is an exact finite sum of
cell measures, computable by a planar arrangement rather than by quadrature. That
makes richer tilings reachable and the answer exact rather than approximate.

\subsection*{Lever 2: tighten $u(29)$ and $u(30)$}
No colouring work at all. Theorem~\ref{thm:B}'s own lower bounds leave room:
$u(29)$ is known only to lie in $[89,103]$ and $u(30)$ in $[93,110]$. Either
$u(29) \le 98$ or $u(30) \le 98$ gives $v_5 \ge 30$ or $v_5 \ge 31$ at once, and
both targets sit inside the intervals that are currently open.

\begin{thebibliography}{9}

\bibitem{AMP}
B. Alexeev, D.~G. Mixon and H. Parshall,
\emph{The Erd\H{o}s unit distance problem for small point sets},
arXiv:2412.11914 (2024, revised 2025).

\bibitem{dGP}
A.~D.~N.~J. de~Grey and J. Parts,
\emph{On lower bounds of the order of $k$-chromatic unit distance graphs},
Geombinatorics \textbf{32} (2022), no.~2; arXiv:2303.14714.

\end{thebibliography}

\end{document}
